\documentclass[a4paper]{article}

\usepackage{graphicx}
\usepackage{amsmath,amssymb}
\usepackage{minted}
\usepackage{multirow}
\usepackage{float}
\usepackage{todonotes}
\usepackage{forest}
\usepackage{xstring}
\usepackage{xcolor}
\usepackage[most]{tcolorbox}
\usepackage{xparse}
\usepackage{natbib}
\usepackage{tikz}

\usetikzlibrary{graphs,graphdrawing}
\usegdlibrary{layered}

% External Graphviz support
\input{/home/donkarlo/Dropbox/repo/nd_ai_project/src/nd_ai/knoweldge_representation/language/formal/computer/programming/tex/latex/code_clips/external_graphviz.tex}

% Semantic tag macros
\input{/home/donkarlo/Dropbox/repo/nd_ai_project/src/nd_ai/knoweldge_representation/language/formal/computer/programming/tex/latex/parts/control_sequence/kind/semtag/semtag.tex}

% Link macros
\input{/home/donkarlo/Dropbox/repo/nd_ai_project/src/nd_ai/knoweldge_representation/language/formal/computer/programming/tex/latex/code_clips/seehere.tex}

\title{Waveform}
\author{Mohammad Rahmani}
\date{}

\begin{document}
    \maketitle
    
    
    \section{FitzHugh--Nagumo model}
        
        The FitzHugh--Nagumo model (FHN)\cite{cebrian-lacasa-2025-six-decades-of-the-fitzhugh-nagumo-model-a-guide-through-its-spatio-temporal-dynamics-and-influence-across-disciplines} describes a prototype of an excitable system, such as a neuron.
        It is a simplified dynamical model of neuronal excitability and action-potential-like behavior.
        
        \begin{equation}
            \dot{v}
            =
            c
            \Bigl(
            v
            -
            \frac{v^{3}}{3}
            +
            w
            +
            I
            \Bigr)
        \end{equation}
        
        \begin{equation}
            \dot{w}
            =
            -
            \frac{1}{c}
            \Bigl(
            v
            -
            a
            +
            bw
            \Bigr)
        \end{equation}
        
        \cite{cebrian-lacasa-2025-six-decades-of-the-fitzhugh-nagumo-model-a-guide-through-its-spatio-temporal-dynamics-and-influence-across-disciplines}
        reviews six decades of research on the FitzHugh--Nagumo model as a simplified but powerful description of excitable dynamics.
        It shows that the model captures not only neuronal excitability and oscillations, but also bistability, traveling waves, synchronization, and spatial pattern formation.
        The paper emphasizes that the model has become useful beyond neuroscience, including cardiac dynamics, biology, electronics, optics, mathematics, and physics.
        Its main contribution is to organize the known dynamical regimes and bifurcations of the model, especially in spatially extended and coupled systems.
    
    
    
            
            \bibliographystyle{plainnat}
            \bibliography{/home/donkarlo/Dropbox/repo/nd_shared_research/bibliography/bibliography.bib}
\end{document}